% What differs between PGF 3.1.11a and 3.1.12, and what LaTeX's own changes reach: forms a profile could be wrong in on one tree only.  DESIGN section 10, "Versions".
\documentclass{article}
\usepackage{amsmath,amssymb}
\usepackage{tikz-cd}
\pdfpagewidth=5in \pdfpageheight=5in
%%%% case bp-old-rule
%%%% pre
\makeatletter
\def\pgf@sys@bp#1{{\pgf@x=#1\relax\pgf@x=0.99627\pgf@x\edef\temp{\expandafter\Pgf@geT\the\pgf@x\space}\pgfutil@toks@\expandafter\expandafter\expandafter{\expandafter\pgfsysprotocol@currentprotocol\temp}\xdef\pgfsysprotocol@currentprotocol{\the\pgfutil@toks@}}}
\makeatother
%%%% body
\begin{tikzcd} A \arrow[r,"f"] & B \end{tikzcd}
%%%% case bp-new-rule
%%%% pre
\makeatletter
\def\pgf@sys@bp@correct#1{#1=\expandafter\pgf@sys@bp@correct@\number#1;400/803;}
\def\pgf@sys@bp@correct@#1#2;#3;{\dimexpr\the\numexpr(2*#1#2\if#1--\else+\fi 1)*#3\relax sp\relax}
\def\pgf@sys@bp#1{{\pgf@x=#1\relax\pgf@sys@bp@correct\pgf@x\edef\temp{\expandafter\Pgf@geT\the\pgf@x\space}\pgfutil@toks@\expandafter\expandafter\expandafter{\expandafter\pgfsysprotocol@currentprotocol\temp}\xdef\pgfsysprotocol@currentprotocol{\the\pgfutil@toks@}}}
\makeatother
%%%% body
\begin{tikzcd} A \arrow[r,"f"] & B \end{tikzcd}
%%%% case bp-correct-changed
%%%% pre
\makeatletter
\def\pgf@sys@bp@correct#1{#1=0.5#1}
\makeatother
%%%% body
\begin{tikzcd} A \arrow[r,"f"] & B \end{tikzcd}
%%%% case bp-correct-at-changed
%%%% pre
\makeatletter
\def\pgf@sys@bp@correct@#1#2;#3;{\dimexpr\the\numexpr(#1#2\if#1--\else+\fi 1)*#3\relax sp\relax}
\makeatother
%%%% body
\begin{tikzcd} A \arrow[r,"f"] & B \end{tikzcd}
%%%% case geT-changed
%%%% pre
\makeatletter
{\catcode`\p=12\catcode`\t=12\gdef\Pgf@geT#1pt{0}}
\makeatother
%%%% body
\begin{tikzcd} A \arrow[r,"f"] & B \end{tikzcd}
%%%% case ltrhbox-changed
%%%% pre
\makeatletter
\def\pgfsys@ltrhbox{\hbox\bgroup\typeout{[[ltr]]}\egroup\hbox}
\makeatother
%%%% body
\begin{tikzcd} A \arrow[r,"f"] & B \end{tikzcd}
%%%% case savedirections-changed
%%%% pre
\makeatletter
\def\pgfsys@savedirections{\typeout{[[save]]}}
\def\pgfsys@restoredirections{\typeout{[[restore]]}}
\makeatother
%%%% body
\begin{tikzcd} A \arrow[r,"f"] & B \end{tikzcd}
%%%% case scope-old
%%%% pre
\makeatletter
\def\pgfscope{\pgfsyssoftpath@setcurrentpath\pgfutil@empty\pgfsys@beginscope\pgf@resetpathsizes\edef\pgfscope@linewidth{\the\pgflinewidth}\let\pgfscope@stroke@color=\pgf@strokecolor@global\let\pgfscope@fill@color=\pgf@fillcolor@global\begingroup}
\makeatother
%%%% body
\begin{tikzcd} A \arrow[r,"f"] & B \end{tikzcd}
%%%% case version-text
%%%% body
\typeout{[[\pgfversion]]}
\begin{tikzcd} A \arrow[r,"f \pgfversion"] & \pgfversiondate \end{tikzcd}
%%%% case dash-numbers
\begin{tikzcd} A \arrow[r,dashed] \arrow[d,dotted] \arrow[rd,dash pattern=on 3.3pt off 1.7pt] \arrow[r,bend left,dash pattern=on 0.1pt off 0.3pt on 7.77pt off 0.07pt] & B \\ C & D \end{tikzcd}
%%%% case dash-phase
\begin{tikzcd} A \arrow[r,dash phase=1pt,dashed] & B \end{tikzcd}
%%%% case widths
\begin{tikzcd} A \arrow[r,line width=0.33pt] \arrow[d,line width=1.7pt] \arrow[rd,ultra thick,bend left=7] & B \\ C \arrow[r,very thin] \arrow[ru,semithick,hook] & D \end{tikzcd}
%%%% case tips
\begin{tikzcd} A \arrow[r,hook] \arrow[d,two heads] \arrow[rd,tail,dashed] & B \arrow[d,maps to] \arrow[r,Rightarrow] & E \arrow[d,harpoon] \\ C \arrow[r,dash] \arrow[ru,no head] & D \arrow[r,leftrightarrow] \arrow[ru,equal] & F \arrow[u,hookleftarrow,bend right] \end{tikzcd}
%%%% case tips-tikz
\begin{tikzcd}[arrow style=tikz] A \arrow[r,hook] \arrow[d,two heads] \arrow[rd,tail,dashed] & B \arrow[d,maps to] \arrow[r,Rightarrow] & E \\ C \arrow[r,dash] & D \arrow[r,leftrightarrow] & F \end{tikzcd}
%%%% case tips-lm
\begin{tikzcd}[arrow style=Latin Modern] A \arrow[r,hook] \arrow[d,two heads] \arrow[rd,tail,dashed] & B \arrow[d,maps to] \arrow[r,Rightarrow] & E \\ C \arrow[r,dash] & D \arrow[r,leftrightarrow] & F \end{tikzcd}
%%%% case loops
\begin{tikzcd} A \arrow[loop above] \arrow[loop below,"x"] \arrow[loop left,dashed] \arrow[loop right,"y"'] & B \arrow[loop,out=30,in=60,distance=7mm] \end{tikzcd}
%%%% case bends
\begin{tikzcd} A \arrow[r,bend left=13] \arrow[r,bend right=47,"f"] \arrow[rd,out=17,in=153,looseness=1.3] & B \\ & C \arrow[lu,to path={-- (\tikztotarget)}] \end{tikzcd}
%%%% case shifts
\begin{tikzcd} A \arrow[r,shift left=0.7] \arrow[r,shift right=1.3] \arrow[rd,xshift=2.3pt,yshift=-1.7pt] \arrow[d,start anchor=south west,end anchor=north east] & B \\ C \arrow[ru,shorten=2.1pt,"f"{pos=0.13}] & D \end{tikzcd}
%%%% case crossing
\begin{tikzcd} A \arrow[rd] & B \arrow[ld,crossing over] \\ C & D \end{tikzcd}
%%%% case phantom
\begin{tikzcd} A \arrow[rd,phantom,"\ulcorner",very near start] & B \\ C \arrow[ru,phantom,"\scriptstyle a"{sloped}] & D \end{tikzcd}
%%%% case big-negative
\begin{tikzcd}[row sep=3.17cm,column sep=4.21cm] A \arrow[r,"f"] \arrow[d] & B \arrow[d,bend left=67] \arrow[ld,"g"'] \\ C \arrow[r] \arrow[ru,bend right=3] & D \arrow[u,bend left=33] \end{tikzcd}
%%%% case colours
%%%% pre
\definecolor{cc}{rgb}{0.13,0.77,0.31}
\colorlet{dd}{cc!37!blue}
%%%% body
\begin{tikzcd} A \arrow[r,red,"f"] \arrow[d,cc,"g",dashed] \arrow[rd,color=dd] & \textcolor{cc}{B} \arrow[d,blue!30] \\ C & D \arrow[l,"h"{text=red},dd!20!red] \end{tikzcd}
%%%% case sci-shift
\begin{tikzcd} A \arrow[r,shift left=1e-1] & B \end{tikzcd}
%%%% case sci-neg-shift
\begin{tikzcd} A \arrow[r,shift left=-5e-1] & B \end{tikzcd}
%%%% case sci-pos
\begin{tikzcd} A \arrow[r,"f"{pos=3e-1}] & B \end{tikzcd}
%%%% case sci-bend
\begin{tikzcd} A \arrow[r,bend left=3e1] & B \end{tikzcd}
%%%% case sci-big-bend
\begin{tikzcd} A \arrow[r,bend left=1.5e1] & B \end{tikzcd}
%%%% case sci-width
\begin{tikzcd} A \arrow[r,line width=1e-1pt] & B \end{tikzcd}
%%%% case sci-sep
\begin{tikzcd}[column sep=1e1pt,row sep=1.5e1pt] A \arrow[r] & B \end{tikzcd}
%%%% case sci-paren
\begin{tikzcd} A \arrow[r,shift left={(1e-1)}] & B \end{tikzcd}
%%%% case sci-shorten
\begin{tikzcd} A \arrow[r,shorten >=2e-1pt] & B \end{tikzcd}
%%%% case sci-distance
\begin{tikzcd} A \arrow[loop,out=3e1,in=6e1,distance=5e0mm] & B \end{tikzcd}
%%%% case reserved-alt-arrow
\begin{tikzcd} A \arrow[r,alt=x] & B \end{tikzcd}
%%%% case reserved-actualtext-arrow
\begin{tikzcd} A \arrow[r,actualtext=x] & B \end{tikzcd}
%%%% case reserved-artifact-arrow
\begin{tikzcd} A \arrow[r,artifact] & B \end{tikzcd}
%%%% case reserved-diagram
\begin{tikzcd}[alt=x] A \arrow[r] & B \end{tikzcd}
%%%% case reserved-cell
\begin{tikzcd} |[alt=x]| A \arrow[r] & B \end{tikzcd}
%%%% case reserved-label
\begin{tikzcd} A \arrow[r,"f"{alt=x}] & B \end{tikzcd}
%%%% case monotone
\begin{tikzcd} A \arrow[r,smooth monotone] & B \end{tikzcd}
%%%% case active-semi-token
%%%% pre
\begingroup\catcode`\;=\active \gdef;{\mathrm{X}}\gdef\semi{;}\endgroup
%%%% body
\begin{tikzcd} \semi \arrow[r] & B \end{tikzcd}
%%%% case active-colon-token
%%%% pre
\begingroup\catcode`\:=\active \gdef:{\mathrm{X}}\gdef\colontok{:}\endgroup
%%%% body
\begin{tikzcd} \colontok \arrow[r] & B \end{tikzcd}
%%%% case active-bar-token
%%%% pre
\begingroup\catcode`\|=\active \gdef|{\mathrm{X}}\gdef\bartok{|}\endgroup
%%%% body
\begin{tikzcd} a\bartok \arrow[r] & B \end{tikzcd}
%%%% case active-bar-token-label
%%%% pre
\begingroup\catcode`\|=\active \gdef|{\mathrm{X}}\gdef\bartok{|}\endgroup
%%%% body
\begin{tikzcd} A \arrow[r,"\bartok"] & B \end{tikzcd}
%%%% case active-semi-label
%%%% pre
\begingroup\catcode`\;=\active \gdef;{\mathrm{X}}\gdef\semi{;}\endgroup
%%%% body
\begin{tikzcd} A \arrow[r,"\semi"] & B \end{tikzcd}
%%%% case active-semi-catcode
%%%% pre
\catcode`\;=\active \def;{\mathrm{X}}
%%%% body
\begin{tikzcd} a; \arrow[r,"f;"] & B \end{tikzcd}
\catcode`\;=12
%%%% case active-colon-catcode
%%%% pre
\catcode`\:=\active \def:{\mathrm{X}}
%%%% body
\begin{tikzcd} a: \arrow[r,"f:"] & B \end{tikzcd}
\catcode`\:=12
%%%% case active-bar-catcode
%%%% pre
\catcode`\|=\active \def|{\mathrm{X}}
%%%% body
\begin{tikzcd} a| \arrow[r,"f|"] & B \end{tikzcd}
\catcode`\|=12
%%%% case active-semi-catcode-other-after
%%%% pre
\begingroup\catcode`\;=\active \gdef;{\mathrm{X}}\gdef\semi{;}\endgroup
\catcode`\;=\active
%%%% body
\begin{tikzcd} \semi; \arrow[r] & B \end{tikzcd}
\catcode`\;=12
%%%% case math-active-colon-undefined
%%%% pre
\mathcode`\:="8000
%%%% body
\begin{tikzcd} A \arrow[r] & B \end{tikzcd}
%%%% case math-active-colon-nouse
%%%% pre
\begingroup\catcode`\:=\active \gdef:{\mathrel{\mathrm{X}}}\endgroup
\mathcode`\:="8000
%%%% body
\begin{tikzcd} a \arrow[r] & B \end{tikzcd}
%%%% case centercolon-nouse
%%%% pre
\usepackage{mathtools}
\mathtoolsset{centercolon}
%%%% body
\begin{tikzcd} a \arrow[r,"x"] & B \end{tikzcd}
$a:b$
%%%% case babel-french-semi
%%%% pre
\usepackage[french]{babel}
%%%% body
\begin{tikzcd} a;b \arrow[r,"x:y;z"] & \text{c; d: e!} \end{tikzcd}
%%%% case babel-french-shorthandoff
%%%% pre
\usepackage[french]{babel}
\AtBeginDocument{\shorthandoff{;:}}
%%%% body
\begin{tikzcd} a;b \arrow[r,"x:y;z"] & \text{c; d: e!} \end{tikzcd}
%%%% case babel-lib
%%%% pre
\usepackage[french]{babel}
\usetikzlibrary{babel}
%%%% body
\begin{tikzcd} a;b \arrow[r,"x:y;z"] & \text{c; d: e!} \end{tikzcd}
%%%% case handle-active-off
%%%% pre
\usepackage[french]{babel}
\tikzset{handle active characters in code=false,handle active characters in nodes=false}
%%%% body
\begin{tikzcd} a;b \arrow[r,"x:y;z"] & \text{c; d: e!} \end{tikzcd}
%%%% case handle-active-on
%%%% pre
\tikzset{handle active characters in code,handle active characters in nodes}
%%%% body
\begin{tikzcd} a;b \arrow[r,"x:y;z"] & \text{c; d: e!} \end{tikzcd}
%%%% case bending
%%%% pre
\usetikzlibrary{bending}
%%%% body
\begin{tikzcd} A \arrow[r,bend left] \arrow[d,bend right=47,"f"] & B \arrow[loop right] \\ C & D \end{tikzcd}
%%%% case bending-meta
%%%% pre
\usetikzlibrary{bending,arrows.meta}
%%%% body
\begin{tikzcd} A \arrow[r,bend left,-{Latex[bend]}] & B \\ C \arrow[r,bend left=60,"f"] & D \end{tikzcd}
%%%% case bending-flex
%%%% pre
\usetikzlibrary{bending}
\tikzset{every picture/.append style={flex}}
%%%% body
\begin{tikzcd} A \arrow[r,bend left=70] & B \end{tikzcd}
%%%% case scope-before
%%%% pre
\makeatletter
%%%% body
\pgfscope \begin{tikzcd} A \arrow[r] & B \end{tikzcd}\endpgfscope
%%%% case group-before
\bgroup \begin{tikzcd} A \arrow[r] & B \end{tikzcd}\egroup
\begingroup \begin{tikzcd} A \arrow[r] & B \end{tikzcd}\endgroup
%%%% case cell-grouptype
\begin{tikzcd} \the\currentgrouptype \arrow[r] & B \end{tikzcd}
%%%% case cell-begingroup
\begin{tikzcd} \begingroup a\endgroup \arrow[r] & \bgroup b\egroup \end{tikzcd}
%%%% case cell-aftergroup
\begin{tikzcd} {\aftergroup x y} \arrow[r] & \begingroup\aftergroup z w\endgroup \end{tikzcd}
%%%% case cell-endgroup-bad
\begin{tikzcd} a\endgroup \arrow[r] & B \end{tikzcd}
%%%% case cell-egroup-bad
\begin{tikzcd} a\egroup\bgroup b \arrow[r] & B \end{tikzcd}
%%%% case hook-env-before
%%%% pre
\AddToHook{env/tikzcd/before}{\typeout{[[before]]}}
%%%% body
\begin{tikzcd} A \arrow[r] & B \end{tikzcd}
%%%% case hook-env-begin
%%%% pre
\AddToHook{env/tikzcd/begin}{\typeout{[[begin]]}\def\zz{z}}
%%%% body
\begin{tikzcd} A\zz \arrow[r] & B \end{tikzcd}
%%%% case hook-env-end
%%%% pre
\AddToHook{env/tikzcd/end}{\typeout{[[end]]}}
%%%% body
\begin{tikzcd} A \arrow[r] & B \end{tikzcd}
%%%% case hook-env-after
%%%% pre
\AddToHook{env/tikzcd/after}{\typeout{[[after]]}}
%%%% body
\begin{tikzcd} A \arrow[r] & B \end{tikzcd}
%%%% case hook-cmd-before
%%%% pre
\AddToHook{cmd/tikzcd/before}{\typeout{[[cmdbefore]]}}
%%%% body
\begin{tikzcd} A \arrow[r] & B \end{tikzcd}
%%%% case hook-cmd-after
%%%% pre
\AddToHook{cmd/tikzcd/after}{\typeout{[[cmdafter]]}}
%%%% body
\begin{tikzcd} A \arrow[r] & B \end{tikzcd}
%%%% case hook-cmd-end-before
%%%% pre
\AddToHook{cmd/endtikzcd/before}{\typeout{[[endbefore]]}}
%%%% body
\begin{tikzcd} A \arrow[r] & B \end{tikzcd}
%%%% case hook-cmd-end-after
%%%% pre
\AddToHook{cmd/endtikzcd/after}{\typeout{[[endafter]]}}
%%%% body
\begin{tikzcd} A \arrow[r] & B \end{tikzcd}
%%%% case hook-para-begin
%%%% pre
\AddToHook{para/begin}{\typeout{[[para begin]]}}
%%%% body
\begin{tikzcd} A \arrow[r] & B \end{tikzcd}

text \begin{tikzcd} A \arrow[r,"\text{a}"] & \text{b} \end{tikzcd}
%%%% case hook-para-end
%%%% pre
\AddToHook{para/end}{\typeout{[[para end]]}}
\AddToHook{para/after}{\typeout{[[para after]]}}
\AddToHook{para/before}{\typeout{[[para before]]}}
%%%% body
\begin{tikzcd} A \arrow[r] & B \end{tikzcd}

text \begin{tikzcd} A \arrow[r,"\text{a}"] & \text{b} \end{tikzcd}
%%%% case hook-selectfont
%%%% pre
\AddToHook{selectfont}{\typeout{[[sf]]}}
%%%% body
\begin{tikzcd} A \arrow[r] & \text{\itshape b} \end{tikzcd}
%%%% case hook-next-env
\AddToHookNext{env/tikzcd/before}{\typeout{[[next]]}}
\begin{tikzcd} A \arrow[r] & B \end{tikzcd}
\begin{tikzcd} A \arrow[r] & B \end{tikzcd}
%%%% case hook-shipout
%%%% pre
\AddToHook{shipout/before}{\typeout{[[ship]]}}
\AddToHook{shipout/background}{\put(0,0){x}}
\AddToHook{shipout/foreground}{\put(0,0){y}}
%%%% body
\begin{tikzcd} A \arrow[r] & B \end{tikzcd}
%%%% case math-hooks
%%%% pre
\makeatletter
\@ifundefined{hook_math}{}{}
\AddToHook{math/before}{\typeout{[[math before]]}}
\AddToHook{math/after}{\typeout{[[math after]]}}
\AddToHook{math/begin}{\typeout{[[math begin]]}}
\AddToHook{math/end}{\typeout{[[math end]]}}
\AddToHook{math/inline/before}{\typeout{[[il before]]}}
\AddToHook{math/inline/after}{\typeout{[[il after]]}}
\AddToHook{math/inline/begin}{\typeout{[[il begin]]}}
\AddToHook{math/inline/end}{\typeout{[[il end]]}}
\makeatother
%%%% body
\begin{tikzcd} A \arrow[r] & B \end{tikzcd}
%%%% case env-copy
%%%% pre
\NewEnvironmentCopy{mycd}{tikzcd}
%%%% body
\begin{mycd} A \arrow[r] & B \end{mycd}
%%%% case cmd-copy
%%%% pre
\NewCommandCopy\mycd\tikzcd
\NewCommandCopy\endmycd\endtikzcd
%%%% body
\begin{mycd} A \arrow[r] & B \end{mycd}
%%%% case env-wrap
%%%% pre
\NewDocumentEnvironment{wcd}{}{\begin{tikzcd}}{\end{tikzcd}}
%%%% body
\begin{wcd} A \arrow[r] & B \end{wcd}
%%%% case env-wrap-body
%%%% pre
\NewDocumentEnvironment{wcd}{+b}{\begin{tikzcd} #1 \end{tikzcd}}{}
%%%% body
\begin{wcd} A \arrow[r] & B \end{wcd}
%%%% case env-renew
%%%% pre
\RenewEnvironment{tikzcd}{\typeout{[[mine]]}}{\typeout{[[end mine]]}}
%%%% body
\begin{tikzcd} A \end{tikzcd}
%%%% case env-new-cmd
%%%% pre
\NewDocumentCommand\mycd{m}{\begin{tikzcd} #1 \end{tikzcd}}
%%%% body
\mycd{A \arrow[r] & B}
%%%% case link-target
%%%% pre
\usepackage{hyperref}
%%%% body
\begin{tikzcd} \phantomsection A \arrow[r] & \MakeLinkTarget[sec]{}B \end{tikzcd}
%%%% case link-target-star
%%%% pre
\usepackage{hyperref}
%%%% body
\begin{tikzcd} \MakeLinkTarget*{foo}A \arrow[r] & B \end{tikzcd} \hyperlink{foo}{x}
%%%% case label-in-cell
%%%% pre
\usepackage{hyperref}
%%%% body
\begin{tikzcd} \phantomsection\label{q}A \arrow[r] & B \end{tikzcd} \ref{q} \pageref{q}
%%%% case label-in-cell-plain
\begin{tikzcd} \label{q}A \arrow[r] & \refstepcounter{equation}\label{r}B \end{tikzcd} \ref{q} \pageref{r} \ref{r}
%%%% case index-in-cell
%%%% pre
\usepackage{imakeidx}
\makeindex
%%%% body
\begin{tikzcd} \index{foo}A \arrow[r,"\index{bar}f"] & B \end{tikzcd}
%%%% case label-index-let
%%%% pre
\usepackage{imakeidx}
\makeindex
%%%% body
\begin{tikzcd} \label{r}A \arrow[r] & B \end{tikzcd} \ref{r}
